% Part: first-order-logic
% Chapter: syntax-and-semantics
% Section: Structures

\documentclass[../../../include/open-logic-section]{subfiles}

\begin{document}

\olfileid{fol}{syn}{str}

\olsection{પ્રથમ-ક્રમ ભાષાઓ માટેની \printtoken{P}{structure}}

\begin{explain}
પ્રથમ-ક્રમ ભાષાઓ પોતે
\emph{અર્થઘટનરહિત} હોય છે: !!{constant}s, !!{function}s અને
!!{predicate}s સાથે કોઈ ચોક્કસ અર્થ જોડાયેલો હોતો નથી. અર્થ
\article{structure} \emph{!!{structure}} નિર્દિષ્ટ કરીને આપવામાં આવે
છે. તે \emph{વ્યાપકક્ષેત્ર}, એટલે કે !!{constant}s જે વસ્તુઓને નિર્દેશે છે,
!!{function}s જે વસ્તુઓ પર કાર્ય કરે છે અને પરિમાણકો જેના પર વ્યાપ્ત
થાય છે તે વસ્તુઓ, નિર્દિષ્ટ કરે છે. તે ઉપરાંત, કયાં !!{constant}s કઈ
વસ્તુઓને નિર્દેશે છે, !!a{function} વસ્તુઓનું વસ્તુઓમાં કેવી રીતે
પ્રતિચિત્રણ કરે છે અને !!{predicate}s કઈ વસ્તુઓને લાગુ પડે છે તે પણ
નિર્દિષ્ટ કરે છે. !!^{structure}s તર્કશાસ્ત્રમાં અર્થવિચારી ખ્યાલો,
જેમ કે ફલિતતા, પ્રામાણ્ય અને સંતોષ્યતાના ખ્યાલોનો પાયો છે. સાહિત્યમાં
તેમને જુદી જુદી રીતે ``સંરચનાઓ'', ``અર્થઘટનો'' અથવા ``નિદર્શો'' કહે
છે.
\end{explain}

\begin{defn}[!!^{structure}s]
પ્રથમ-ક્રમ તર્કશાસ્ત્રની ભાષા~$\Lang{L}$ માટેની \Article{structure}
\emph{!!{structure}}~$\Struct M$ નીચેના ઘટકોની બનેલી હોય છે:
\begin{enumerate}
\item \emph{વ્યાપકક્ષેત્ર:} અરિક્ત ગણ~$\Domain M$
\item \emph{!!{constant}sનું અર્થઘટન:} $\Lang{L}$ના દરેક
  !!{constant}~$c$ માટે !!a{element}~$\Assign{c}{M} \in \Domain M$
\item \emph{!!{predicate}sનું અર્થઘટન:} $\Lang{L}$ના દરેક $n$-સ્થાની
  !!{predicate}~$R$ માટે ($\eq$ સિવાય) $n$-સ્થાની સંબંધ
  $\Assign{R}{M} \subseteq \Domain{M}^n$
\item \emph{!!{function}sનું અર્થઘટન:} $\Lang{L}$ના દરેક $n$-સ્થાની
  !!{function}~$f$ માટે $n$-સ્થાની વિધેય
  $\Assign{f}{M} \colon \Domain{M}^n \to \Domain{M}$
\end{enumerate}
\end{defn}

\begin{ex}
અંકગણિતની ભાષા માટેની !!^a{structure}~$\Struct M$માં એક ગણ,
!!{constant}~$\Obj 0$ના અર્થઘટન તરીકે $\Domain M$નો એક ઘટક
$\Assign{\Obj 0}{M}$, એકસ્થાની વિધેય
$\Assign{\Obj \prime}{M} \colon \Domain{M} \to \Domain M$, બે
દ્વિસ્થાની વિધેયો $\Assign{\Obj +}{M}$ અને $\Assign{\Obj
  \times}{M}$ (બંને $\Domain M^2 \to \Domain M$) અને દ્વિસ્થાની
સંબંધ $\Assign{\Obj <}{M} \subseteq \Domain{M}^2$ હોય છે.

આવી સંરચનાનો એક સ્વાભાવિક દાખલો આ છે:
\begin{enumerate}
\item $\Domain N = \Nat$
\item $\Assign{\Obj 0}{N} = 0$
\item દરેક $n \in \Nat$ માટે $\Assign{\Obj \prime}{N}(n) = n + 1$
\item દરેક $n, m \in \Nat$ માટે $\Assign{\Obj +}{N}(n, m) = n + m$
\item દરેક $n, m \in \Nat$ માટે $\Assign{\Obj \times}{N}(n, m) = n\cdot m$
\item $\Assign{\Obj <}{N} = \Setabs{\tuple{n, m}}{n \in \Nat, m \in
  \Nat, n < m}$
\end{enumerate}
આ રીતે વ્યાખ્યાયિત~$\Lang L_A$ માટેની સંરચના~$\Struct N$ને
\emph{અંકગણિતનો પ્રમાણભૂત નિદર્શ} કહે છે, કારણ કે તે~$\Lang L_A$ના
અતાર્કિક અચળોનું આપણે અપેક્ષા રાખીએ એ જ રીતે અર્થઘટન કરે છે.

તેમ છતાં,~$\Lang L_A$ માટે બીજી ઘણી !!{structure}s શક્ય છે. દાખલા
તરીકે, વ્યાપકક્ષેત્ર તરીકે~$\Nat$ને બદલે પૂર્ણાંકોનો ગણ~$\Int$ લઈ શકીએ અને
$\Obj 0$, $\Obj \prime$, $\Obj +$, $\Obj \times$, $\Obj <$નાં
અર્થઘટનો તે મુજબ વ્યાખ્યાયિત કરી શકીએ. પરંતુ~$\Lang L_A$ માટે એવી
સંરચનાઓ પણ વ્યાખ્યાયિત કરી શકીએ જેને સંખ્યાઓ સાથે દૂરનો પણ સંબંધ ન
હોય.
\end{ex}

\begin{ex}
ગણસિદ્ધાંતની ભાષા~$\Lang L_Z$ માટેની સંરચના~$\Struct M$ને માત્ર એક
ગણ અને એક દ્વિસ્થાની સંબંધની જરૂર પડે છે. તેથી, તાંત્રિક રીતે, લોકોનો
ગણ અને ``$x$,~$y$ કરતાં મોટી ઉંમરનું છે'' એ સંબંધ, અથવા~$\Nat$ સાથે
$n, m \in \Nat$ માટે $n \ge m$,~$\Lang L_Z$ માટે !!a{structure}
તરીકે વાપરી શકાય.
\sourcecorrection{OLSEM-001}{મૂળની \texttt{structures.tex},
પંક્તિ~91માં \texttt{single-two place relation} એવો વિકૃત સંયુક્ત
પ્રયોગ છે. ભાષાને માત્ર એક દ્વિસ્થાની સંબંધ જોઈએ છે; તેથી અહીં
\texttt{single two-place relation}નો અર્થ પુનઃસ્થાપિત કર્યો છે.}

$\Lang L_Z$ માટેની એક ખાસ રસપ્રદ !!{structure}, જેમાં વ્યાપકક્ષેત્રના
!!{element}s ખરેખર ગણો છે અને $\Obj \in$નું અર્થઘટન ખરેખર ``$x$ એ
~$y$નો !!a{element} છે'' એ સંબંધ છે, તે \emph{આનુવંશિક રીતે સાન્ત
ગણો}ની !!{structure}~$\Struct{{HF}}$ છે:
\begin{enumerate}
\item $\Domain{{{HF}}} = \emptyset \cup \Pow{\emptyset} \cup
  \Pow{\Pow{\emptyset}} \cup \Pow{\Pow{\Pow{\emptyset}}} \cup \dots$;
\item $\Assign{\Obj \in}{{{HF}}} = \Setabs{\tuple{x, y}}{x, y \in
  \Domain{{{HF}}}, x \in y}$.
\end{enumerate}
\end{ex}

\begin{digress}
!!a{structure} તરીકે શું ગણાય તે વિશેની આપણી શરતો આપણા તર્કશાસ્ત્રને
અસર કરે છે. દાખલા તરીકે, ખાલી વ્યાપકક્ષેત્રો ન સ્વીકારવાની પસંદગી એ સુનિશ્ચિત
કરે છે કે પરિમાણિત વાક્યો માટેના સામાન્ય સંતોષવિવરણ (અથવા
સત્યતાવિવરણ) હેઠળ $\lexists[x][(!A(x) \lor \lnot !A(x))]$
પ્રમાણભૂત---એટલે કે તાર્કિક સત્ય---છે. અને બધાં !!{constant}s વ્યાપકક્ષેત્રની
કોઈ વસ્તુને નિર્દેશે જ એવી શરત અસ્તિત્વવાચક વ્યાપકીકરણને અનુમાનનો
સુદૃઢ ઢાંચો બનાવે છે: $!A(a)$, તેથી $\lexists[x][!A(x)]$. જો આપણે
નામોને વ્યાપકક્ષેત્ર બહારની વસ્તુ નિર્દેશવા અથવા કશું ન નિર્દેશવા દઈએ, તો
આપણે એવા \emph{મુક્ત તર્કશાસ્ત્ર} તરફ જઈએ જેમાં અસ્તિત્વવાચક
વ્યાપકીકરણ માટે વધારાનું આધારવિધાન જોઈએ: $!A(a)$ અને
$\lexists[x][\eq[x][a]]$, તેથી $\lexists[x][!A(x)]$.
\end{digress}

\end{document}
